\chapter{今後の方針}

本章では，現時点までに得られた成果をまとめ，卒業論文の完成に向けた今後の方針を述べる．

\section{現時点のまとめ}

本研究は，日本語タスクにおいて確信度の引き出し方が LLM の較正に与える影響を明らかにすることを目的としている．本論文の時点で得られた成果は以下の通りである．

第一に，3 つの確信度引き出し方式を設計し，これらを同一条件で比較する評価基盤を実装した．基盤は ECE，Brier Score，AUROC，Murphy 分解を計算し，問題の復元抽出による信頼区間を算出する．応答全文を保存する設計としたため，採点処理の不具合が後から見つかった際に，追加の問い合わせなしに再採点できた．

第二に，日本語版 MGSM の 250 問を用いた予備実験を行った．Ling.1S が他の2 方式より ECE，Brier Score，AUROC，正答率において劣ることが，信頼区間によって支持された．一方，REL と RES では差を検出できず，Verb.1S と Verb.2S の差はどの指標でも判定できなかった．

第三に，3 方式に共通して確信度が 0.9 以上に偏ることを観察した．この分布のもとでは ECE の値が平均確信度と平均正答率の差に近づき，較正の良否と正答率の高低を読み分けにくいことを示した．当初は Murphy 分解の REL によってこの読み分けが可能と考えていたが，REL も各区間の正答率を含むため分離の手段にはならず，実測でも REL による差の検出はできなかった．

第四に，方法論上の知見として，採点処理の詳細が方式間の比較に影響しうることを確認した．回答の抽出に関わる不具合が方式によって偏って発生しており，修正によって正答率の差が 12〜15 ポイントから 8〜10 ポイントへ縮小した．

\section{今後の計画}

\subsection{採点処理の検証と保存仕様の改善}

最優先の作業は採点処理の信頼性を高めることである．現在までに修正した不具合は外部のレビューで指摘された事例に基づくものであり，網羅的な検証は行っていない．応答の形式を分類し，抽出と判定が意図通りかを標本調査で確認する．このとき，正解を見てから抽出規則を調整することは避ける．規則を先に定め，その結果を評価する順序を守る．

あわせて，Verb.2S の第 1 ターンの応答全文を保存するよう実装を改める．現在の保存仕様では，Verb.2S だけ回答の再抽出ができず，他の 2 方式と検証可能性が非対称になっている．

\subsection{較正のみを取り出す比較条件}

6.2 節で述べた，較正の違いと正答率の違いを分離できないという問題に対処する．すべての方式に共通の回答をあらかじめ与え，その回答に対する確信度のみを尋ねる条件を設ける．正誤ラベルが方式間で共通になるため，較正の違いだけを比較できる．

この条件は，第 1 章に挙げた第 2 の問いに直接答えるものであり，指標の選択を工夫するより確実である．

\subsection{確信度を尋ねない基準条件}

6.3 節で述べた，Ling.1S の正答率低下の原因を調べるために，確信度を尋ねずに回答のみを求める条件を設ける．これにより，各方式が正答率をどの程度変化させるかを測れる．

なお，この条件は Verb.2S の第 1 ターンと実質的に同じである．7.2.1 で保存仕様を改めれば，Verb.2S の実行結果から同時に得られる．ただし現在保存されている結果からは取り出せないため，改めて実行する必要がある．

\subsection{確信度の偏りを対象としたプロンプト設計}

本研究の中心課題は，較正を改善するプロンプトの設計である．予備実験により，確信度が上端に偏り，方式を変えてもほとんど動かないことが分かった．

検討している方向は 3 つある．

\textbf{誤りうる理由の列挙を先行させる方向．} 確信度を答える前に，その回答が誤っている可能性のある理由を挙げさせる．自信過剰が，誤りの可能性を検討しないまま確信度を述べることに由来するのであれば，検討を促すことで確信度が変化すると期待される．

\textbf{確信度の使用分布に制約を与える方向．} Ling.1S において最上位に偏っていたことを踏まえ，各段階を使う頻度に関する指示をプロンプトに加える．ただしこの方向は，較正を改善したのではなく分布を操作しただけになる恐れがあり，評価の解釈に注意を要する．

\textbf{問題間の相対比較を求める方向．} 複数の問題を提示し，確信の高い順に並べさせる．絶対的な確率の推定より順序の判断のほうが容易であれば，識別の面で改善が期待できる．ただしこの条件は，提示する問題の組み合わせや順序の影響を受けるうえ，順位を確率に変換する規則を定めなければ ECE や Brier Scoreを計算できない．評価基盤の変更も要るため，他の 2 方向より実施の負担が大きい．

ここで注意すべきは，\textbf{分布の偏りが解消することを改善の条件としない}ことである．正答率が高い課題では高い確信度に偏ること自体は妥当でありうる．また分布を広げるだけでは較正は改善しない．採否は較正，識別，正答率，そして費用を合わせて判断する．

\subsection{データセットの拡張}

MGSM の対象が算数の文章題に限られるという制約に対処するため，JCommonsenseQA と JMMLU による追試を行う．3.3.3 項で述べたとおり，JCommonsenseQA は選択肢が 5 つであるため，記号の生成と照合の実装を改める必要がある．JMMLU については，翻訳による科目と日本固有の内容を扱う科目のどちらを用いるかを明示する．

問題数が増えることで，5.4.4 項で判定できなかった Verb.1S と Verb.2S の比較についても検出力が上がる．ただし別のデータセットでは効果の大きさ自体が変わりうるため，標本数を増やせば決着するとは限らない．

\subsection{対象モデルの拡大}

現在は claude-sonnet-4-6 のみを対象としている．他の商用モデルおよび日本語を主対象としたモデルを加え，観察された傾向がモデルに依存しないかを確認する．とくに確信度の上端への偏りがモデル共通の性質であるかは，本研究の主張の一般性に関わる．

\subsection{言語表現の対応づけに関する感度の確認}

6.4 節で限界として挙げた，言語表現と数値の対応づけの暫定性に対処する．割り当てを変えた複数の条件で指標を再計算し，Ling.1S の評価がどの程度対応づけに依存するかを調べる．ECE，Brier Score，REL は変化するが，段階の順序と同順位の関係を保つ対応づけであれば AUROC は変化しない．この作業は保存済みの応答から再計算できるため，追加の問い合わせを必要としない．

\section{作業計画}

今後の作業の順序を表 7.1 に示す．

\begin{table}[htbp]
\centering
\small
\caption{今後の作業計画}
\label{tab:7.1}
\begin{tabular}{|l|l|l|l|}
\hline
順序 & 作業 & 追加の実行 & 対応する節 \\
\hline
1 & 採点処理の検証と保存仕様の改善 & 不要 & 7.2.1 \\
2 & 言語表現の対応づけの感度確認 & 不要 & 7.2.7 \\
3 & 較正のみを取り出す比較条件の実行 & 要 & 7.2.2 \\
4 & 確信度を尋ねない基準条件の実行 & 要 & 7.2.3 \\
5 & プロンプト条件の設計と少数条件での検証 & 要 & 7.2.4 \\
6 & データセットを分けたうえでの評価 & 要 & 7.2.5 \\
7 & 対象モデルの拡大 & 要 & 7.2.6 \\
\hline
\end{tabular}
\end{table}

順序 1 と 2 を先に置くのは，追加の問い合わせを必要とせず，かつ以降のすべての結果の信頼性に関わるためである．採点が正しいことを確かめないまま条件を増やすと，誤りを含んだまま規模だけが大きくなる．

順序 5 と 6 については，条件を選ぶために用いるデータと，最終的な評価に用いるデータを分ける．同じデータを見ながら条件を選び，そのデータで最終評価を行うと，選択の影響が評価に入り込む．最終評価の前に，比較する条件と主要な指標を確定させる．
